\subsection{Prekladanie DNA/mRNA do signálu}
Prekladanie DNA alebo RNA sekvencií do elektrického signálu v nanopórovom sekvenovaní \cite{Zheng2023Nanopore} zahŕňa niekoľko kľúčových krokov:
\begin{itemize}
    \item \textbf{Príprava vzorky:} DNA alebo RNA molekuly sú pripravené na sekvenovanie, často zahŕňajúce fragmentáciu a pridanie adaptérnych sekvencií.
    \item \textbf{Vedenie cez nanopór:} Molekuly sú vedené cez nanopór pomocou elektrického poľa, ktoré spôsobuje ich pohyb \cite{Zheng2023Nanopore}.
    \item \textbf{Detekcia prúdu:} Keď molekula prechádza cez nanopór, každá báza spôsobuje charakteristickú zmenu v elektrickom prúde, ktorý je zaznamenaný ako časová séria dát \cite{Zheng2023Nanopore}.
    \item \textbf{Analýza signálu:} Zaznamenaný elektrický signál je analyzovaný pomocou algoritmov na identifikáciu sekvencie nukleotidov na základe zmien prúdu \cite{Zheng2023Nanopore}.
    \item \textbf{Preklad do sekvencie:} Softvér prekladá analyzovaný signál späť do sekvencie DNA alebo RNA, čo umožňuje výskumníkom získať genetické informácie z pôvodnej molekuly \cite{Zheng2023Nanopore}.
\end{itemize}
